\chapter{Glucagon}

\section*{What Is This Test?}
Glucagon is a 29-amino-acid single-chain polypeptide hormone produced by the alpha cells of the pancreatic islets of Langerhans. Glucagon is the primary counter-regulatory hormone to insulin, raising blood glucose through stimulation of hepatic glycogenolysis (breakdown of glycogen to glucose) and gluconeogenesis (synthesis of glucose from non-carbohydrate precursors). Glucagon secretion is stimulated by hypoglycemia, amino acids (especially arginine and alanine), sympathetic nervous system activation (via beta-adrenergic receptors), and cholecystokinin/gastrin. It is inhibited by hyperglycemia, insulin, somatostatin, amylin, and free fatty acids. Glucagon is synthesized as proglucagon and is also produced in the intestinal L-cells (where it is cleaved to GLP-1, GLP-2, oxyntomodulin, and glicentin rather than glucagon). Glucagon is rapidly cleared by the liver and kidneys, with a plasma half-life of 3--6 minutes. The glucagon blood test is rarely ordered clinically but has specific indications: diagnosis of glucagon-secreting tumors (glucagonomas, rare pancreatic neuroendocrine tumors), evaluation of unexplained hypoglycemia, and research applications. More commonly, glucagon is used as a pharmacologic agent: in glucagon stimulation testing for GH deficiency or pheochromocytoma evaluation, and in the emergency treatment of severe hypoglycemia.

\section*{Alternative Names}
\begin{itemize}
  \item Serum Glucagon
  \item Plasma Glucagon
  \item Pancreatic Glucagon
  \item Glucagonoma Marker
  \item Glucagon-Like Peptide-1 (GLP-1, distinct test)
  \item Enteroglucagon (gut-derived, distinct from pancreatic glucagon)
\end{itemize}
\section*{Principle}
Glucagon is measured by competitive radioimmunoassay (RIA) or enzyme-linked immunosorbent assay (ELISA). Because of glucagon's low circulating concentration (~25--50 pg/mL) and rapid degradation, the assay requires a highly specific antibody against the C-terminal region of glucagon (amino acids 19--29) to distinguish pancreatic glucagon from gut-derived enteroglucagon and proglucagon fragments (glicentin, oxyntomodulin, GLP-1, GLP-2), which share sequence homology. Aprotinin (a protease inhibitor) is added to the collection tube to prevent glucagon degradation by plasma proteases. Blood must be collected into a pre-chilled EDTA tube containing aprotinin, placed on ice immediately, centrifuged at 4 degrees C within 30 minutes, and plasma frozen at -20 degrees C or colder. Even with these precautions, glucagon is one of the most unstable hormones in clinical measurement, and results from improperly handled specimens are unreliable. Most hospital laboratories do not perform glucagon measurement in-house; it is a send-out test to specialized reference laboratories. Mass spectrometry-based methods for glucagon have been developed but are primarily research tools \cite{tietz2023textbook}.

\section*{History}
Glucagon was discovered in 1923 by Charles Kimball and John Murlin at the University of Rochester, who noted that pancreatic extracts initially caused a transient rise in blood glucose (glucagon effect) before the more prolonged fall (insulin effect). The amino acid sequence of glucagon was determined by William Bromer and colleagues in 1957. The first radioimmunoassay for glucagon was developed by Roger Unger and colleagues in 1959 at UT Southwestern. The syndrome of glucagonoma (necrolytic migratory erythema, diabetes, weight loss, glossitis, and elevated glucagon) was described by McGavran, Unger, and colleagues in 1966. The distinction between pancreatic glucagon and gut-derived glucagon-like immunoreactivity (enteroglucagon) was established in the 1970s, leading to the development of pancreas-specific glucagon assays. GLP-1 was discovered in the 1980s as a cleavage product of proglucagon in intestinal L-cells and was subsequently developed into a major therapeutic class (GLP-1 receptor agonists: exenatide, liraglutide, semaglutide) for type 2 diabetes.

\section*{Why This Test Is Ordered}
\begin{itemize}
  \item Diagnosis of glucagonoma --- a rare pancreatic neuroendocrine tumor (incidence 1 in 20 million per year) causing the glucagonoma syndrome: necrolytic migratory erythema (characteristic skin rash --- erythematous, blistering, crusted lesions in groin, perineum, extremities that migrate and heal with hyperpigmentation), diabetes mellitus or glucose intolerance, weight loss, deep vein thrombosis (glucagon increases factor X production), glossitis, cheilitis, diarrhea, and neuropsychiatric symptoms. Fasting glucagon is typically >500 pg/mL (often >1000 pg/mL) \cite{wermers1996glucagonoma}.
  \item Investigation of suspected multiple endocrine neoplasia type 1 (MEN1) --- glucagonoma can be part of MEN1 syndrome
  \item Differential diagnosis of hypoglycemia --- inappropriately elevated glucagon during hypoglycemia suggests dysfunctional counter-regulation; absent glucagon response during hypoglycemia is characteristic of type 1 diabetes (loss of alpha cell glucagon secretion, contributing to hypoglycemia unawareness)
  \item Research evaluation of alpha cell function in diabetes mellitus
  \item Rarely, diagnosis of glucagon receptor defects causing pancreatic alpha cell hyperplasia (Mahvash disease) with marked hyperglucagonemia, pancreatic neuroendocrine tumors, and mild diabetes
  \item Evaluation of familial glucagonoma syndromes
\end{itemize}

\section*{Primary vs Secondary Test}
Glucagon is a rare, specialized test ordered for the specific indication of suspected glucagonoma. It is not a primary or routine clinical test. The diagnosis of glucagonoma is typically confirmed by markedly elevated fasting glucagon (>500--1000 pg/mL) with clinical features and pancreatic mass on imaging. In hypoglycemia evaluation, insulin, C-peptide, and proinsulin are the primary tests; glucagon may be a secondary adjunct in specialized research or refractory cases.

\section*{How This Test Is Performed}
A healthcare professional draws a fasting blood sample into a pre-chilled EDTA (purple-top) tube containing aprotinin (a protease inhibitor). The tube must be placed on ice immediately, centrifuged at 4 degrees C within 30 minutes, and the plasma aliquoted and frozen at -20 degrees C or colder. Improper specimen handling (room temperature, delayed processing, no aprotinin) results in glucagon degradation and falsely low values. This is a send-out test processed at specialized reference laboratories.

\section*{What Preparation Is Needed}
The patient must fast for 10--12 hours overnight before testing because meals stimulate glucagon secretion. Strenuous exercise should be avoided for 24 hours. The specimen requires the specialized handling described above. Medications affecting glucagon should be noted: GLP-1 receptor agonists (exenatide, liraglutide, semaglutide) suppress glucagon; DPP-4 inhibitors (sitagliptin, linagliptin) increase GLP-1 and suppress glucagon; insulin suppresses glucagon; somatostatin analogues (octreotide, lanreotide) suppress glucagon; and amino acid infusions stimulate glucagon. High-dose biotin should be held for 48 hours.

\section*{Contraindications and Precautions}
\begin{itemize}
  \item The most critical pre-analytical factor is specimen collection and handling: glucagon is extremely unstable in vitro. Specimens collected without aprotinin or not processed on ice within 30 minutes will yield falsely low results that are clinically misleading. Even with optimal handling, glucagon values are assay-dependent and reference ranges vary considerably. Glucagon levels must be interpreted in the context of simultaneous glucose, insulin, and C-peptide levels, and clinical symptoms. Stress (acute illness, surgery, trauma) can elevate glucagon 2--4 fold. Hypoglycemia is the primary physiologic stimulus for glucagon; an elevated glucagon at the time of hypoglycemia may be a normal counter-regulatory response, not indicative of a tumor. Glucagon elevation without clinical features of glucagonoma syndrome is common in: diabetes mellitus (relative or absolute insulin deficiency eliminates tonic inhibition of alpha cells), chronic kidney disease (decreased renal clearance of glucagon), liver cirrhosis (decreased hepatic glucagon clearance), acute pancreatitis, starvation, and severe stress. In glucagonoma, the plasma glucagon level typically far exceeds (>500--1000 pg/mL) any of these non-tumor elevations (<200--300 pg/mL).
\end{itemize}
\section*{Risks}
Minimal risk from venipuncture. Slight pain, bruising, or hematoma. Rare: vasovagal syncope. No risks specific to glucagon testing beyond standard venipuncture.

\section*{What Happens After the Test}
Results are typically available within 7--14 days (send-out test). Fasting glucagon >500--1000 pg/mL with clinical features of glucagonoma syndrome (necrolytic migratory erythema, weight loss, diabetes) confirms the diagnosis \cite{eldor2011glucagonoma}. The next step is localization: CT or MRI of the abdomen (glucagonomas are often large, >3--5 cm, located in the pancreatic tail, hypervascular on contrast-enhanced CT), followed by somatostatin receptor imaging (Ga-68 DOTATATE PET/CT). Surgical resection is the treatment of choice. If glucagon is mildly elevated (<200 pg/mL) without the characteristic syndrome, other causes (diabetes, renal failure, cirrhosis, stress) should be considered, and the result may not be clinically significant. In a patient with hypoglycemia: if glucagon is appropriately elevated (>100 pg/mL during hypoglycemia), the counter-regulatory response is intact. If glucagon fails to rise during hypoglycemia, there is alpha cell dysfunction, most commonly from type 1 diabetes (loss of intra-islet insulin-mediated alpha cell regulation).

\section*{Understanding Results}

\subsection*{Below Normal}
\begin{itemize}
  \item Type 1 diabetes mellitus: loss of the glucagon response to hypoglycemia is an early and profound defect in type 1 diabetes, occurring within the first few years of disease. Absent glucagon secretion combined with loss of epinephrine response (autonomic neuropathy) causes severe hypoglycemia unawareness. This is due to loss of beta cells and the paracrine insulin signal that normally regulates alpha cells, not destruction of alpha cells themselves \cite{unger2012glucagonocentric}.
  \item Pancreatectomy: removal of pancreatic alpha cells eliminates glucagon production. Patients are at risk for severe hypoglycemia due to combined insulin and glucagon deficiency.
  \item Somatostatin analogue therapy (octreotide, lanreotide): somatostatin suppresses glucagon secretion from alpha cells. Used therapeutically to control symptoms in glucagonoma and other functional neuroendocrine tumors.
  \item GLP-1 receptor agonist therapy (exenatide, liraglutide, semaglutide): GLP-1 suppresses glucagon secretion as part of its glucose-lowering mechanism. This contributes to their therapeutic efficacy.
  \item Insulin-induced hypoglycemia: exogenous insulin administration causing hypoglycemia suppresses endogenous glucagon secretion (a normal response to hyperinsulinemia).
\end{itemize}

\subsection*{Above Normal}
\begin{itemize}
  \item Glucagonoma: markedly elevated fasting glucagon (>500 pg/mL, often >1000--5000 pg/mL). Classic glucagonoma syndrome includes necrolytic migratory erythema (present in 70--80\%, pathognomonic skin rash that waxes and wanes), weight loss (90\%), diabetes mellitus or glucose intolerance (75--80\%), diarrhea (30\%), venous thromboembolism (30\%, glucagon increases factor X production), glossitis/stomatitis/cheilitis, neuropsychiatric symptoms (depression), and normochromic normocytic anemia. Tumor is usually solitary, >3--5 cm, in the pancreatic tail. Approximately 50--80\% have metastases at presentation. Treatment: surgical resection, somatostatin analogues for symptom control, and chemotherapy or targeted therapy for metastatic disease \cite{vanbeek2004glucagonoma}.
  \item Diabetes mellitus: relative or absolute insulin deficiency eliminates tonic inhibition of alpha cells, causing inappropriate hyperglucagonemia that contributes to fasting and postprandial hyperglycemia. Glucagon typically 100--200 pg/mL (much lower than glucagonoma) \cite{dunning2007alpha}.
  \item Chronic kidney disease (CKD stages 4--5): decreased renal clearance and degradation of glucagon, with levels 2--3 fold elevated.
  \item Liver cirrhosis: decreased hepatic glucagon clearance. May be associated with glucose intolerance.
  \item Physiologic response to hypoglycemia: glucagon is rapidly secreted in response to hypoglycemia (<70 mg/dL) to stimulate hepatic glycogenolysis and gluconeogenesis. Plasma glucagon rises within minutes to 100--300 pg/mL. This is a normal counter-regulatory response.
  \item Acute stress: trauma, burns, sepsis, surgery, and myocardial infarction elevate glucagon as part of the hypermetabolic stress response (along with cortisol and catecholamines).
  \item Acute pancreatitis: glucagon may be elevated due to islet cell damage and inflammation.
  \item Mahvash disease (glucagon cell hyperplasia and neoplasia): inactivating glucagon receptor mutations cause marked hyperglucagonemia (glucagon receptors on alpha cells mediate negative feedback). Pancreatic alpha cell hyperplasia and neuroendocrine tumors develop. Mild diabetes or normal glucose tolerance because the receptor defect prevents hyperglycemia. Extremely rare.
  \item Familial glucagonoma: associated with MEN1 syndrome. Glucagonoma is the third most common functional pancreatic NET in MEN1 after gastrinoma and insulinoma.
\end{itemize}

\section*{Normal Results}
\begin{itemize}
  \item Fasting plasma glucagon (adults): 25--50 pg/mL (lab-specific, with significant variability)
  \item Fasting plasma glucagon (reference labs): 20--100 pg/mL (depending on assay specificity for pancreatic vs. total glucagon immunoreactivity)
  \item Glucagon in glucagonoma: >500 pg/mL (usually >1000 pg/mL)
  \item Glucagon response to hypoglycemia: increases 3--10 fold above baseline to 100--300 pg/mL
  \item Glucagon response to intravenous arginine (stimulation test): increases 3--10 fold
  \item Note: Glucagon assays are poorly standardized and values vary significantly between laboratories. Results must be interpreted in the context of clinical findings, glucose, insulin, and C-peptide levels. Always use the reference range from the performing laboratory.
\end{itemize}

\section*{Follow-up Tests}
\begin{itemize}
  \item Fasting glucose and insulin/C-peptide: Essential co-tests. Glucagon excess in glucagonoma causes glucose intolerance or diabetes (insulin resistance and increased hepatic glucose output). In hypoglycemia evaluation, glucose, insulin, and C-peptide are primary.
  \item Chromogranin A: General neuroendocrine tumor marker that is often elevated in glucagonoma and other pancreatic NETs. Used for monitoring disease progression.
  \item Pancreatic polypeptide, VIP, gastrin, somatostatin: Other pancreatic hormones to screen for in suspected MEN1 or multi-hormonal NET.
  \item CT or MRI of the abdomen (pancreatic protocol, triple-phase): Glucagonomas are typically hypervascular, well-defined masses >3--5 cm, most commonly in the pancreatic tail. CT angiography phase detects hypervascular NETs.
  \item Ga-68 DOTATATE PET/CT: Somatostatin receptor-based imaging with superior sensitivity for detecting primary and metastatic glucagonoma (most express somatostatin receptors, particularly sst2). Now the preferred functional imaging modality.
  \item Endoscopic ultrasound (EUS) with FNA: For small pancreatic NETs not well visualized on CT/MRI.
  \item Genetic testing for MEN1: In patients with glucagonoma and family history or other MEN1 features (hyperparathyroidism, pituitary adenoma, other pancreatic NETs).
  \item Prothrombin time, D-dimer: Screen for venous thromboembolism (VTE risk elevated in glucagonoma due to factor X induction).
  \item Skin biopsy: Necrolytic migratory erythema has characteristic histology (necrolytic keratinocytes in upper epidermis).
  \item Zinc and amino acid levels: Glucagonoma-induced hypoaminoacidemia contributes to the skin rash and responds to amino acid and zinc supplementation while awaiting surgery.
\end{itemize}

\section*{References}
\bibliographystyle{unsrtnat}
\bibliography{references}

\clearpage
\section*{Regulatory Status and Authorization Date}
This chapter describes a test category, not a specific commercial product. FDA, CDSCO, EU IVDR, and MHRA approval or registration dates therefore cannot be assigned without the assay manufacturer, product name, instrument, and intended use. For laboratory-developed tests, an individual agency approval date may not exist. See the Preface for the interpretation of regulatory status.

\clearpage
